\author{Peter McHale}
\documentclass{article}
\usepackage{amsmath, amssymb}
\usepackage{microtype}
\usepackage{xcolor}
\title{How bias arises in Gnocchi}
\begin{document}
\maketitle

\section{Expected number of SNVs under neutrality}
In a selectively neutral genomic window,
the expected number of sites that are polymorphic 
(a site is polymorphic if an ALT allele is seen 
at that site in one or more members of a given cohort)
is
\begin{equation}
    \sum_{i} p_{c_i}(x_i)
    \label{eq:expected_snvs}
\end{equation}
where:
\begin{itemize}
    \item $i$ indexes the single-nucleotide sites in the window
    \item $x_i$ is a vector of regional features (GC content, etc)
    associated with site $i$,
    \item $c_i$ is the sequence and methylation context of the site $i$,
    \item $p_c(x)$ is the probability that a site is polymorphic,
    conditional on context $c$ and regional feature vector $x$,
    and, of course, the number of individuals in the cohort.
\end{itemize}

\section{Assuming uniformity of the features within a genomic window}

Assume that all the features vectors $\{x_i\}$, ranging over all sites $i$ in a window $w$,
are approximately the same: 
\[ x_i = x_w. \]
Under this assumption, equation~\eqref{eq:expected_snvs} simplifies to
\begin{equation}
\sum_{i} p_{c_i}(x_w) = \sum_c n_c^{(w)} p_c(x_w),
\end{equation}
where the right side groups probabilities that appear on the left side by 
context, and $n_c^{(w)}$ is the number of sites in window $w$ that have 
sequence and methylation context $c$. 

\section{The naive way to fit a model of the polymorphism probability}

To predict the SNV count of a window with feature vector $x$, we need to 
model $p_c(x)$, the probability that a site with context $c$ 
and feature vector $x$ is polymorphic (as opposed to monomorphic) across 
the cohort of interest. 
The natural approach is to 
construct a parametric model of $p_c(x)$ as follows: 
\begin{equation} 
    p_c(x) = \sigma(\beta_c \cdot x),
    \label{eq:joint_model}
\end{equation}
where: 
\begin{itemize}
    \item $\sigma(x)$ is the sigmoid function (bounded between 0 and 1), 
    \item $\cdot$ is the dot product,
    \item $\beta_c$ is a parameter vector associated with the sequence and methylation context $c$,
\end{itemize}
and then train one model for each context $c$ on an appropriate genome-wide 
training set of sites 
(each with context $c$, but with varying feature vectors $x$) 
that are either monomorphic or polymorphic in the cohort. 

\section{The training set is contaminated with selection that may correlate with $x$}

The problem with this approach lies in choosing the training set of sites. 
These sites should be selectively neutral, but, 
in practice we don't know how to identify such sites in the noncoding genome, 
where functional annotations are infrequent and imprecise. 
Moreover, noncoding sites under selection may correlate with $x$
(e.g., promoters, which are under selection, tend to have high GC content), causing $p_c(x)$ 
to be inaccurate in the tails of the $x$-distribution. 

\section{A 2-stage approach to eliminate the effect of feature-dependent contamination of the training set}

De novo mutations (DNMs), by definition, have not
been exposed to natural selection, and therefore
inform us about the inherent mutability of a genomic site. 
Can we use DNMs as a training set? 
Not immediately: we want a neutral 
estimate of the polymorphism probability, which scales with cohort size, 
not DNM rate, which does not. 

There is a way, however, to draw that DNM rate 
into the polymorphism probability such that it eliminates the effect of 
feature-dependent contamination with selection that would otherwise occur in a training set 
of polymorphisms constructed from the entire noncoding genome. 
To leverage a set of neutrally pure DNMs, 
Chen et al broke the problem into two steps: 
\begin{equation}
    p_c(x) = r_c(x) \, p_c(\bar{x}),
\end{equation}
where $\bar{x}$ is the average of $x$ over all windows of interest, 
and
\begin{equation}
    r_c(x) = \frac{p_c(x)}{p_c(\bar{x})},
\end{equation}
is a feature-dependent multiplicative adjustment to the feature-independent 
$p_c(\bar{x})$.

In the first step, they approximated $p_c(\bar{x})$
by the maximum-likelihood estimate: 
\begin{equation}
    p_c(\bar{x}) \approx \frac{N_c^{(p)}}{N_c^{(m)} + N_c^{(p)}}
\end{equation}
where 
\begin{itemize} 
    \item $N_c^{(m)}$ is the number of sites in the noncoding genome with 
    context $c$ that are monomorphic, and
    \item $N_c^{(p)}$ is the number of sites in the noncoding genome with 
    context $c$ that are polymorphic
\end{itemize}
Given the large number of training examples, and the fact that 
most of them have feature vectors close to $\bar{x}$, 
this estimate is likely to be quite accurate
(though subject to low levels of contamination by selection).

In the second step, they assumed that the polymorphism probability
in the cohort, $p_c(x)$, 
is approximately proportional to the DNM probability, $\hat{p}_c(x)$,
in an arbitrary training set of DNMs and background sites (not DNMs).
Under this assumption:
\begin{equation}
    r_c(x) \approx \frac{\hat{p}_c(x)}{\hat{p}_c(\bar{x})}
\end{equation}
where 
\begin{equation}
\hat{p}_c(x) = \sigma(\hat{\beta}_c \cdot x),
\end{equation}
with parameters $\hat{\beta}_c$ estimated 
from the training set of DNMs and background sites (not DNMs).
Notice that, even though $\hat{p}_c(x)$ is sensitive to the class balance in the training set
(the ratio of DNMs examples to non-DNM examples), 
the estimate of $r_c(x)$, which is a ratio of such probabilities, is not.
Though this training set can be trusted to be neutral, even in the tails of the $x$-distributon, 
it is much smaller than that used to compute $p_c(\bar{x})$, 
so the estimate of $r_c(x)$ is likely to be noisier than that of $p_c(\bar{x})$
solely because of sampling noise.  

\section{Hypothesis: the feature-dependent adjustment is biased}
As we show in our manuscript, 
the process of training any function of $x$ works, in general, to minimize 
global bias at the expense of local bias in the tails of the $x$-distribution. 
Applied to the adjustment factor $r_c(x)$, whose training data is sparse,
we expect $r_c(x)$ to be accurate near $\bar{x}$, but inaccurate (and possibly biased) 
in the tails ($x$ far from $\bar{x}$). 
This bias would propagate, via $p_c(x)$, to the expected per-window SNV counts as a function of $x$,
potentially explaining the bias we observe in Gnocchi at extreme GC content, for example.  

\section{Predictions of the hypothesis}
It is instructive to consider what would happen, under the above hypothesis, 
as the size of the DNM training set is increased. 

When the DNM training set is so small that the tails are not sampled at all, 
then $r_c(x)$ is expected to tend towards its value at $\bar{x}$, 
which is 1 by construction, 
and Gnocchi reduces to a context-only model, 
inheriting whatever bias it has as a function of $x$. 
\textcolor{red}{We performed this
experiment and found
not only that bias in Gnocchi was reduced, but
also that it approached the levels of bias we saw
in other context-only constraint metrics.}

As the DNM training set size increases, and more and more extreme 
values of $x$ are sampled, 
$r_c(x)$ will start to capture the true function near $\bar{x}$
making Gnocchi better there, 
but this may come at the expense of deviating from the true function away 
from $\bar{x}$, in a way that we cannot predict without knowing the truth. 
\textcolor{red}{This is what we believe is causing the 
large bias we observe in Gnocchi.}

Eventually, if enough DNM training data existed to densely sample extreme $x$, 
then we would expect the bias at extreme $x$ to disappear. 
\textcolor{red}{We approximated this regime by increasing the number of 
background 
sites (only) in the DNM training set, 
retrained $r_c(x)$, and observed that bias below 
the level we reported in our paper.}

Together these experiments lend support to our hypothesis that 
the sparseness of the DNM training set is responsible for the large bias 
observed in Gnocchi as a function of GC content. 

\section{Derivation of the quantities plotted in Figure 3}
\label{sec:fig3}

Figure~3 plots three quantities, all binned on the GC content of a 1\,kb window.
This section derives each of them from the model above, and states exactly which
approximations enter and which do not.

\subsection{Notation}

Write the context $c = (t, m)$, where $t$ is the trinucleotide and $m$ the
methylation level. This distinction matters here, and is not cosmetic: step~1
conditions on both $t$ and $m$, whereas the fitted adjustment of step~2 conditions
on $t$ only. Chen et al.\ fit one model per trinucleotide, pooling all methylation
levels.

Let $W$ be the analyzed set of $N$ windows (noncoding, passing quality control,
autosomes and pseudoautosomal regions), let $g$ index a GC bin, and let
$W_g \subset W$ be the windows falling in bin $g$. For a window $w$, let $O(w)$ be
the observed number of polymorphic sites in the cohort, and $n_c^{(w)}$ the number
of callable sites of context $c$ in $w$, as before.

Two further points of fact about the fitted model, both confirmed against Chen et
al.'s code rather than their text. First, the operative adjustment is
\begin{equation}
    r_t(x) = \frac{\sigma\!\left(\beta_{t,0} + \beta_t \cdot z(x)\right)}
                  {\sigma\!\left(\beta_{t,0}\right)},
    \label{eq:r_operative}
\end{equation}
a ratio of two \emph{predicted probabilities} from the same fitted model, with
$z(x)$ the standardized, PCA-rotated feature vector and $z(\bar{x}) = 0$ by
construction. (Their Methods state a ratio of linear predictors instead; the code
computes Eq.~\eqref{eq:r_operative}.) Second, and consequentially, $r_t$ is a
\emph{ratio}: any error in $\hat{p}_t$ that is a pure level shift, common to
numerator and denominator, cancels exactly and never reaches the constraint score.
Only feature-driven \emph{variation} in $\hat{p}_t$ survives. Everything below is
organized around that fact.

\subsection{The two expected counts}

The context-only (``step~1'') expected count of a window is
Eq.~\eqref{eq:expected_snvs} with the adjustment set to unity:
\begin{equation}
    E_1(w) = \sum_c n_c^{(w)}\, p_c(\bar{x}),
    \qquad
    E_1^{t}(w) = \sum_{c\,:\,t(c)=t} n_c^{(w)}\, p_c(\bar{x}),
    \label{eq:E1}
\end{equation}
so that $E_1 = \sum_t E_1^t$. The adjusted (``step~2'', published Gnocchi)
expected count applies $r$ per trinucleotide:
\begin{equation}
    E_2(w) = \sum_t E_1^{t}(w)\, r_t(x_w).
    \label{eq:E2}
\end{equation}
Both are available per window and per $(w,t)$ from the pipeline's own exports, so
neither is reconstructed here.

\subsection{Panel A: the mean standardized rank}

For a model $M \in \{1, 2\}$ Chen et al.\ summarize a window by the signed
chi statistic
$z_M(w) = -\operatorname{sign}(O - E_M)\,\sqrt{(O-E_M)^2/E_M}$,
which is just
\begin{equation}
    z_M(w) = \frac{E_M(w) - O(w)}{\sqrt{E_M(w)}},
    \label{eq:z}
\end{equation}
positive when a window carries less variation than expected. Windows are retained
only if $z_M \in [-10,10]$ for every $M$ plotted, so all curves describe one
identical window set. Each $z_M$ is then converted to a standardized rank within
that set,
\begin{equation}
    u_M(w) = \frac{\operatorname{rank}_M(w) - \tfrac{1}{2}}{N} \in (0,1),
\end{equation}
which is uniform on $(0,1)$ by construction, with mean exactly $1/2$. Panel~A
plots the conditional mean
\begin{equation}
    \bar{u}_M(g) = \frac{1}{|W_g|}\sum_{w\in W_g} u_M(w),
    \label{eq:panelA}
\end{equation}
with error bars $\operatorname{sd}(u_M)/\sqrt{|W_g|}$. Because the marginal
distribution of $u_M$ is uniform regardless of the model, $\bar{u}_M(g) = 1/2$ for
every $g$ is exactly the statement that the model has no GC-dependent bias, and
departures from $1/2$ measure that bias on a scale that is comparable across
models with entirely different units. This is also why the third curve, depletion
rank, can be overlaid despite being defined on a different (Halldorsson) window
set: it is ranked within itself, and the sign convention is opposite, so its
complement $1 - \mathrm{DR}$ is plotted.

\subsection{Why panel A is so sensitive to the adjustment}

It is worth recording how strongly Eq.~\eqref{eq:panelA} responds to $r$, because
this is what licenses reading panel~B as the cause of panel~A. Suppose a window's
adjustment is a uniform inflation, $E_2 = f E_1$, and write the observed count as
$O = E_1(1+\epsilon)$. Then from Eq.~\eqref{eq:z},
\begin{equation}
    z_2 - z_1 = \sqrt{E_1}\left[\frac{f - 1 - \epsilon}{\sqrt{f}} + \epsilon\right]
              \;\approx\; \sqrt{E_1}\,(f-1)
    \qquad (f \to 1,\ \epsilon \to 0).
    \label{eq:sensitivity}
\end{equation}
The $\sqrt{E_1}$ prefactor is the point. A typical analyzed window has
$E_1 \approx 174$, so $\sqrt{E_1} \approx 13$, and a mere $10\%$ inflation
displaces $z$ by more than one unit --- comparable to the spread of $z$ itself.
Measured directly by re-ranking the real genome-wide $z$ distribution under a
uniform inflation $f$, the mean rank moves from $0.500$ at $f=1$ to $0.687$ at
$f=1.10$ and $0.302$ at $f=0.90$. Small multiplicative errors in $r$ are therefore
not a second-order concern; they are the whole effect.

\subsection{Panel B, first form: the adjustment Gnocchi applies}

Dividing Eq.~\eqref{eq:E2} by Eq.~\eqref{eq:E1} gives the single multiplicative
factor each window actually receives,
\begin{equation}
    \bar{r}(w) = \frac{E_2(w)}{E_1(w)}
               = \sum_t \frac{E_1^{t}(w)}{E_1(w)}\, r_t(x_w),
\end{equation}
an $E_1$-weighted mean of the per-trinucleotide adjustments. Aggregating to a GC
bin as a ratio of sums rather than a mean of ratios,
\begin{equation}
    R(g) = \frac{\sum_{w \in W_g} E_2(w)}{\sum_{w \in W_g} E_1(w)},
    \label{eq:Rg}
\end{equation}
which is the adjustment the bin receives in aggregate, and which keeps the
following decomposition exact.

Let $K = \{\mathrm{ACG},\mathrm{CCG},\mathrm{GCG},\mathrm{TCG}\}$ be the CpG
trinucleotides. Define the CpG share of the step-1 expectation and the two
group-wise adjustments,
\begin{equation}
    \Pi(g) = \frac{\sum_{w\in W_g}\sum_{t\in K} E_1^{t}(w)}{\sum_{w\in W_g} E_1(w)},
    \qquad
    R_{\mathrm{CpG}}(g) = \frac{\sum_{w\in W_g}\sum_{t\in K} E_1^{t}(w)\, r_t(x_w)}
                    {\sum_{w\in W_g}\sum_{t\in K} E_1^{t}(w)},
\end{equation}
and $R_{\mathrm{non}}$ likewise over $t \notin K$. Splitting the numerator and
denominator of Eq.~\eqref{eq:Rg} at $K$ gives, with no approximation whatsoever,
\begin{equation}
    R(g) = \Pi(g)\, R_{\mathrm{CpG}}(g) + \bigl(1 - \Pi(g)\bigr) R_{\mathrm{non}}(g).
    \label{eq:decomposition}
\end{equation}
The first form of panel~B plots the three curves of
Eq.~\eqref{eq:decomposition} against $g$, together with the counterfactual
obtained by holding the non-CpG adjustment at unity,
\begin{equation}
    R^{\mathrm{cf}}(g) = \Pi(g)\, R_{\mathrm{CpG}}(g) + \bigl(1 - \Pi(g)\bigr).
\end{equation}
Empirically $R_{\mathrm{CpG}}(g)$ is flat at $0.99$--$1.00$ across the entire GC range while
$R_{\mathrm{non}}(g)$ climbs from $0.96$ to $1.79$, and $R^{\mathrm{cf}}(g)$ is flat to
within $0.6\%$ even though $\Pi$ reaches $0.43$ at the highest GC. Since
Eq.~\eqref{eq:decomposition} is an identity, this is not a fit but a bookkeeping
statement: the GC dependence of what Gnocchi applies is carried entirely by the
non-CpG contexts.

That CpG models are inert is a direct consequence of the ratio structure noted
after Eq.~\eqref{eq:r_operative}. Chen et al.\ deliberately strip GC content, CpG
island status, and three correlated features from the CpG-context models; what
remains barely varies with GC, so the numerator and denominator of
Eq.~\eqref{eq:r_operative} move together and $r_t \approx 1$. Their large level
miscalibration --- these models cannot represent methylation, and hypomethylated
CpG islands in GC-rich sequence mutate at roughly a third the rate of methylated
CpGs --- cancels, and never reaches the score. Methylation is thus \emph{ruled
out} as a cause of the bias, not implicated in it.

\subsection{Panel B, second form: the adjustment the DNMs support}

Establishing that the non-CpG adjustment drives the bias does not establish that
it is \emph{wrong}. For that we need an empirical counterpart of $r_t$ computed
from data the model was meant to describe. Since $r_t$ is by definition the factor
by which the true mutation rate departs from its context-average value, the
empirical target is a rate, measurable per trinucleotide and GC bin:
\begin{equation}
    \tilde{p}_t(g) = \frac{D_t(g)}{A_t(g)},
    \label{eq:emp_rate}
\end{equation}
where $D_t(g)$ is the number of observed de novo mutations of context $t$ falling
in $W_g$, and $A_t(g)$ is the number of \emph{opportunities} --- possible SNVs of
context $t$ --- in the same windows. Numerator and denominator are restricted to
the same window set, and DNMs are assigned to windows by position, so the
denominator counts exactly the territory the numerator was counted over.

\paragraph{Normalization is forced, not chosen.} Chen et al.'s own assumption
(Sec.~6) is that the polymorphism and DNM probabilities are proportional,
$\hat{p}_t(x) = \lambda_t\, p_t(x)$, with $\lambda_t$ unknown and
\emph{context-specific}: it absorbs both the cohort-size scaling and the fact that
$p_t(\bar{x})$ saturates by different amounts in different contexts. The
proportionality constant cancels in a within-context ratio,
\begin{equation}
    \frac{\hat{p}_t(x)}{\hat{p}_t(\bar{x})}
    = \frac{\lambda_t\, p_t(x)}{\lambda_t\, p_t(\bar{x})}
    = r_t(x),
\end{equation}
and only there. So the empirical adjustment must be formed as the same ratio,
each context normalized by its own $E_1$-weighted average over bins,
\begin{equation}
    \hat{r}^{\mathrm{obs}}_t(g) = \frac{\tilde{p}_t(g)}{\langle \tilde{p}_t\rangle},
    \qquad
    \langle \tilde{p}_t\rangle
      = \frac{\sum_{g'} E_1^{t}(g')\, \tilde{p}_t(g')}{\sum_{g'} E_1^{t}(g')},
    \label{eq:emp_norm}
\end{equation}
where $E_1^t(g) = \sum_{w \in W_g} E_1^t(w)$. No free constant is introduced and
none is fitted; $\lambda_t$ never has to be estimated.

The model side must then be normalized \emph{identically},
$\hat{r}^{\mathrm{mod}}_t(g) = r_t(g) / \langle r_t \rangle$ with
$r_t(g) = E_2^{t}(g)/E_1^{t}(g)$. This is not symmetry for its own sake. As GC
increases the trinucleotide composition of windows shifts; if only the empirical
side were normalized, the residual between-context level differences on the model
side would be reweighted bin by bin and would appear in the aggregate as GC
dependence that no model actually asserts.

The two curves of the second form of panel~B are then the same $E_1$-weighted
aggregate over non-CpG contexts,
\begin{equation}
    R^{\bullet}_{\mathrm{non}}(g)
      = \frac{\sum_{t \notin K} E_1^{t}(g)\, \hat{r}^{\bullet}_t(g)}
             {\sum_{t \notin K} E_1^{t}(g)},
    \qquad \bullet \in \{\mathrm{mod}, \mathrm{obs}\},
\end{equation}
and the lower sub-panel plots their ratio, the over-adjustment
\begin{equation}
    \Phi(g) = \frac{R^{\mathrm{mod}}_{\mathrm{non}}(g)}{R^{\mathrm{obs}}_{\mathrm{non}}(g)}.
    \label{eq:inflation}
\end{equation}
A bin's contribution is used only if it carries enough observed DNMs to estimate
Eq.~\eqref{eq:emp_rate}, and a bin is reported only if the usable contexts still
account for at least $90\%$ of its non-CpG $E_1$ weight, so that no bin is built
from an unrepresentative minority of contexts. The error bar is Poisson on
$D_t(g)$, propagated through Eqs.~\eqref{eq:emp_norm}--\eqref{eq:inflation}; the
model curve carries none, being a deterministic function of already-fixed feature
values.

Measured this way, $R^{\mathrm{mod}}_{\mathrm{non}}$ climbs monotonically from $0.95$
to $1.55$ while $R^{\mathrm{obs}}_{\mathrm{non}}$ stays near unity until
$\mathrm{GC}\approx0.55$, so $\Phi$ crosses $1$ near $\mathrm{GC}=0.40$ and
reaches $1.22$--$1.26$, many standard errors from unity. By
Eq.~\eqref{eq:sensitivity} an over-adjustment of that size is more than enough to
produce the rank displacement in panel~A.

\paragraph{Why opportunities, and not expected counts, in Eq.~\eqref{eq:emp_rate}.}
A rate needs a count of opportunities in its denominator. It happens that the
step-1 expected count would serve equally well here, because a non-CpG
trinucleotide has a single methylation level and therefore a single set of fitted
per-site rates:
\begin{equation}
    E_1^{t}(g) = \sum_{\mathrm{alt}} A_t^{\mathrm{alt}}(g)\, \pi_t^{\mathrm{alt}}
               = A_t(g)\, \mu_t ,
    \qquad
    \mu_t = \tfrac{1}{3}\textstyle\sum_{\mathrm{alt}} \pi_t^{\mathrm{alt}},
\end{equation}
with $\pi_t^{\mathrm{alt}} = p_c(\bar{x})$ the fitted per-site rate, and where the
second equality uses $A_t^{\mathrm{alt}}(g) = A_t(g)/3$: every callable position
of context $t$ contributes one opportunity per alternative allele, so each of the
three alleles carries the same site count. The
per-context constant $\mu_t$ is absorbed by the normalization
Eq.~\eqref{eq:emp_norm}, so the two denominators give the same curve --- verified
to agree within $0.1\%$ at every bin. Replacing $A_t(g)$ instead by a real sample
of matched background sites from the DNM training set, which removes the fitted
model from the denominator altogether, changes the curve by $1$--$2\%$. The result
does not rest on the choice.

\paragraph{The one assumption that does bite: callability.} $D_t(g)$ counts DNMs
anywhere in a window, whereas $A_t(g)$ counts only positions gnomAD can call. That
fraction is not flat in GC --- short-read coverage degrades in GC-rich sequence,
and the callable fraction $\gamma(g)$ falls from $0.91$ at $\mathrm{GC}=0.30$ to
$0.75$ at $\mathrm{GC}=0.61$. If trio DNM ascertainment tracks gnomAD's, the two
are matched and no correction applies. If instead the trio call sets are closer to
complete, the matched denominator is $A_t(g)/\gamma(g)$, so the corrected
empirical rate is $\gamma(g)\,\tilde{p}_t(g)$: since $\gamma$ decreases with GC,
the corrected empirical curve falls further in the tail and $\Phi$ \emph{grows},
from $1.22$ to $1.44$ at $\mathrm{GC}=0.61$. The finding survives either way and
the correction only strengthens it, so the honest statement is a range,
$\Phi \in [1.22, 1.44]$, rather than a point estimate. The residual, untestable
possibility runs the other way: if trio DNM calling is itself \emph{less}
sensitive in GC-rich sequence than gnomAD is, part of the measured
over-adjustment is technical rather than a property of the model.

\section{A proposal for Gnocchi 2.0}
Though DNMs ameliorate the effect of selection contamination of the training set, 
the trade-off is that the sparseness of DNMs introduces bias into the constraint metric, 
precisely where it was hoped DNMs would help. 
Perhaps the quest to avoid contamination of the training set with sites under selection 
was misguided,
and it is better to tolerate some contamination if it yields a less biased constraint score 
where it matters most. 
This suggests an approach to a possibly improved constraint score: 
capture the $x$ dependence in Eq.~(\ref{eq:joint_model}) by training directly on dense gnomAD data, 
instead of factorizing it and training one of the factors on sparse DNM data. 
\textcolor{red}{Since this requires a new training process, we deemed this experiment out-of-scope for this rebuttal.} 

\end{document}
